\chapter{From the Base Repository to the Pakistan-Adapted Thesis System}

\section{Starting Point: the Original Base Repository}

The folder \texttt{Reference Documents and Folders/cyclesgym base} provides the original repository snapshot from which the present thesis work evolved. That base repository already exposed the central idea of CyclesGym: a reinforcement-learning wrapper around the CYCLES crop simulator. It contained environment definitions, training entrypoints, notebooks, and documentation, but it did not yet encode the full thesis-specific layer required for a Pakistan-grounded, defense-ready research artifact.

In the base form, the repository was useful as a simulator interface and experimentation framework. However, its structure was still closer to a research code release than to a thesis platform. Important thesis needs were either missing or underdeveloped: localized price profiles, explicit scope mapping against the proposal, broader nutrient support, standardized per-run summary outputs, and narrative documentation that could tie engineering changes to thesis claims.

\section{What Changed in the Thesis Repository}

The current repository is materially larger and more structured than the base version. New directories such as \texttt{final\_docs}, \texttt{Changes}, \texttt{Experimentation and Results}, \texttt{wandb exports}, and the current thesis workspace itself reflect a transition from a generic research repo to a traceable thesis codebase. The new structure does not merely add notes. It adds engineering layers that are directly relevant to the thesis: localized pricing utilities, thesis reporting callbacks, a hierarchical environment, and a broad matrix runner for the latest Pakistan NPK design.

Figure~\ref{fig:repo-evolution} summarizes this transition at the capability level.

\begin{figure}[H]
    \centering
    \includegraphics[width=0.82\textwidth]{base_to_current_evolution.pdf}
    \caption{Capability evolution from the original base repository to the current thesis stack. The progression is defined by completed engineering phases, not by narrative aspiration alone.}
    \label{fig:repo-evolution}
\end{figure}

The most important lesson from this evolution is that thesis readiness is not equivalent to simply running more experiments. It also requires building the surrounding apparatus that makes those experiments interpretable, reproducible, and bounded in scope.

\section{Architecture of the Current Thesis Stack}

At a high level, the current system has four layers. The top layer is the policy-learning layer, implemented through Stable-Baselines3. The second layer is the environment layer, where \texttt{Corn}, crop-planning environments, and the hierarchical environment map actions into simulator-relevant operations. The third layer is the simulation I/O layer, which writes control, operation, weather, and related files and invokes the CYCLES executable. The fourth layer is the parsing and reporting layer, which reconstructs observations, computes rewards, and writes thesis-grade outputs.

\begin{figure}[H]
    \centering
    \includegraphics[width=0.92\textwidth]{system_architecture_diagram.pdf}
    \caption{Component-level architecture of the current thesis stack. The diagram highlights the interaction between policies, environments, simulator interfaces, and thesis reporting artifacts.}
    \label{fig:system-architecture}
\end{figure}

The architecture matters because the thesis is not just about model performance. It is also about whether the stack is structured in a way that permits clear causal interpretation of results. If the simulator run, output parsing, or reporting layer is opaque, then even strong returns become difficult to defend. The current architecture improves that situation by making the system more modular and by attaching thesis-focused reporting outputs to the same environment steps used for training.

\input{tables/generated/module_responsibility_map.tex}

Table~\ref{tab:module-map} turns the architecture narrative into a file-level map. This is important for thesis defense because the committee can ask not only what the system does, but where the relevant logic lives. A traceable answer is much stronger than an abstract architectural description.

\section{Simulation Lifecycle and Execution Semantics}

The thesis repository still depends on the same fundamental simulation loop: the environment writes inputs, invokes CYCLES, reads outputs, and returns observations and rewards. But the lifecycle has been hardened and documented more carefully than in the base repo, partly because integration issues and cross-platform execution paths had to be stabilized before the broader NPK experiment campaign could be treated seriously.

\begin{figure}[H]
    \centering
    \includegraphics[width=0.90\textwidth]{simulation_sequence_diagram.pdf}
    \caption{Sequence of one environment interaction cycle, from reset and input-file creation to simulator execution, parsing, and reward calculation.}
    \label{fig:simulation-sequence}
\end{figure}

This lifecycle is central to reproducibility. A thesis committee may reasonably ask whether a reported return is really linked to the intended combination of weather file, action implementer, and reward function. Because the system rewrites and re-parses simulator artifacts on every relevant step, the answer must be expressed as a concrete file-and-callback pipeline rather than as a vague software description.

\subsection{Class Responsibilities and Separation of Concerns}

The current repository separates concerns more clearly than the inherited base version. \texttt{Observers} convert simulator outputs into learning-state features. \texttt{Implementers} translate agent actions into operations-file edits. \texttt{Rewarders} calculate scalar economic signals from harvest and nutrient costs. \texttt{Managers} handle file-backed input and output data structures. This separation is not merely stylistic. It made it possible to add Pakistan price localization and thesis reporting without rewriting the entire environment stack.

\section{Five Completed Implementation Phases}

A large share of the thesis engineering effort is captured in the local change logs. These logs are not cited as external literature; instead, they function as internal traceability artifacts that justify why the thesis can describe particular capabilities as implemented. Table~\ref{tab:implementation-phases} condenses those records into a compact narrative.

\input{tables/generated/implementation_phases.tex}

The table reveals an important pattern. The work progressed from realism to dimensionality to integration to reporting to run readiness. This is a more defensible trajectory than jumping directly into a large experiment matrix while the platform still lacks localized economics, multi-nutrient control, or structured reporting outputs.

\subsection{Why the Change Logs Matter as Internal Evidence}

The thesis uses the change logs in the \texttt{Changes/} directory as internal evidence of implementation chronology. They are not treated as external scholarly sources, but they serve an important methodological function: they preserve the order in which capabilities were added and reduce the risk of retroactively narrating the repository as if all features had existed from the beginning. This kind of engineering chronology is particularly valuable in research codebases where the final state can obscure the path by which stability was reached.

\section{Localization Layers Added by the Thesis}

The first localization layer is agronomic timing. The repository now includes a Pakistan crop-calendar utility that supports conservative sowing windows based on official or semi-official guidance. The second localization layer is economic. The price utilities no longer rely only on legacy assumptions; they now load yearly Pakistan-relevant crop and nutrient price series. The third localization layer is narrative. The repo now explicitly documents what it does and does not optimize under the Pakistan framing.

These layers are conceptually separate. Crop calendars are about agronomic plausibility. Price series are about economic validity. Scope documentation is about scientific honesty. They should not be conflated, because each answers a different thesis-risk question.

Another localized design choice is the decision to keep the working experiments on the maize-soy configuration while still exposing crop-calendar hooks that can support broader future crop sets. This preserves forward compatibility without forcing the thesis to defend crop configurations that have not yet been stabilized in code and experiments.

\section{From Nitrogen-Centric Logic to NPK-Aware Logic}

One of the most important thesis-specific upgrades is the shift from a mostly nitrogen-centric pipeline to an NPK-aware pipeline. In the earlier state, many experiments and reward interpretations were naturally read as nitrogen optimization. In the current state, the fertilization environment, pricing utilities, and reporting layer can handle N, P, and K as separate channels. This is scientifically useful even before the full broad matrix is rerun, because it changes the modeling capacity of the system and the meaning of cost decomposition.

At the code level, this upgrade appears in multiple places at once. The \texttt{Corn} environment can now expose a \texttt{MultiDiscrete} NPK action space. The pricing utilities load nutrient-wise Pakistan series rather than a single scalar nitrogen cost. The reward layer computes cost channel by channel. The hierarchical reporting callback records \texttt{n\_kg}, \texttt{p\_kg}, \texttt{k\_kg}, and their corresponding costs separately. The scientific importance of the NPK extension therefore lies not in one isolated code change, but in the way several modules now agree on the same nutrient vocabulary.

This upgrade does not automatically imply that the final empirical chapter can already claim broad NPK superiority. The full fresh campaign is still needed for that. But it does mean that the thesis can legitimately describe the current stack as NPK-capable rather than merely N-only.

\section{Hierarchical Integration as a Thesis-Level Extension}

The addition of the hierarchical environment is significant because it aligns the software representation more closely with how real farm management decisions are made. Some decisions, such as crop choice and planting windows, are yearly or seasonal. Others, such as nutrient application, are repeated within a season. A flat environment can model both kinds of decisions, but a hierarchical structure makes the distinction explicit and easier to report.

That distinction is valuable in two ways. First, it supports clearer reasoning about what the policy is actually doing. Second, it creates a cleaner mapping between the technical stack and the story told in the thesis. A committee can understand annual planning and weekly fertilization as different but related decision layers. The code now reflects that intuition instead of leaving it implicit.

The hierarchical environment also improves thesis reporting. Its \texttt{info} dictionary now carries yearly crop-choice metadata, planting-window compliance, and nutrient cost breakdowns alongside the standard step outputs. This is important because it allows later analysis to separate ``what crop was chosen'' from ``how much nutrient was applied'' instead of compressing both into one final reward number.

\section{Why Engineering Traceability Matters for the Thesis}

A common weakness in experimental software theses is that the final written document refers to a polished conceptual system that the code never fully implemented. This thesis attempts to avoid that pattern by treating change logs, test outcomes, and generated tables as first-class evidence. The fact that the full test suite reached 59 passing tests before the fresh broad campaign is not a decorative software statistic. It is part of the argument that the platform is stable enough to support a defensible experiment phase.

The next chapter builds on this system view by formalizing the methodological choices encoded in the current stack, including state, action, reward, matrix structure, and reporting design.
